% !TEX program = xelatex
% 编译方式：把本文件夹整包复制到纯 ASCII 路径后，执行两遍 xelatex（详见 docs/_manual_common/README.md）。
% 首版骨架：内容取自 src/network_reconfiguration/ 与 include/hacdcpf/network_reconfiguration/ 源码
%           （topology_reconfiguration.hpp / topology_analysis.hpp），公式与参数以代码为准。
\documentclass[UTF8,fontset=windows,a4paper,11pt]{ctexart}
\usepackage{hysim_manual}

\title{\textbf{交直流混合高弹性能源电力系统仿真平台}\\[0.5em]
       {\Large 网络重构技术手册（交直流 MILP 与配电 ONR）}\\[0.3em]
       {\large 数学模型、接口参数与验证校核（首版骨架）}}
\author{HySim-XJTU-HRPES（西安交通大学 HRPES 团队）}
\date{\today}

\begin{document}
\maketitle
\begin{abstract}
\noindent 本手册依据 HySim-XJTU-HRPES 平台网络重构模块
（\srcpath{src/network\_reconfiguration/}、
\srcpath{include/hacdcpf/network\_reconfiguration/}）整理，覆盖两条路径：面向交直流混合系统
的故障后重构 MILP（\srcpath{run\_topology\_reconfiguration}）与面向配电网的最优网络重构
ONR（\srcpath{solve\_optimal\_reconfiguration}，LinDistFlow MILP + 验证潮流）。全部结论以源码与
\srcpath{docs/network\_reconfiguration\_models.md} 为准：LinDistFlow 为线性近似，
\fld{reconf\_loss\_mw} 为拓扑代理量而非潮流损耗；ONR 为 AC-only。两处公共符号位于命名空间
\texttt{hacdcpf::analysis}。文档版式与《碳流追踪与碳排放核算技术手册》一致。
\end{abstract}

\tableofcontents
\newpage

\section{阅读指南与约定}
\subsection{数据来源}
\begin{itemize}
  \item \srcpath{include/hacdcpf/network\_reconfiguration/topology\_reconfiguration.hpp}、
        \srcpath{src/network\_reconfiguration/topology\_reconfiguration.cpp} —— 交直流混合重构 MILP；
  \item \srcpath{include/hacdcpf/network\_reconfiguration/topology\_analysis.hpp}、
        \srcpath{src/network\_reconfiguration/topology\_analysis.cpp} —— 配电 ONR 与拓扑查询；
  \item \srcpath{docs/network\_reconfiguration\_models.md} —— 既有模型契约。
\end{itemize}

\subsection{单位制与索引空间}
支路以 \texttt{BranchRef}（\fld{category}$\in\{$\texttt{AC\_Line},\texttt{DC\_Line},
\texttt{VSC\_Coupling}$\}$ 加 \fld{index}）稳定标识，避免混合系统中 AC/DC/VSC 的 \fld{index}
重叠含糊；电压以标幺值（pu）且 LinDistFlow 中 $v_i=|V_i|^2$，功率 MW/MVA，损耗代理
\fld{reconf\_loss\_mw} 单位 MW，时限 \fld{max\_time\_s} 单位 s。ONR 旧式整数 ID 向量
（\fld{open\_branch\_ids} 等）仅在纯 AC 语境无歧义。

\subsection{诚实口径}
LinDistFlow 为线性近似（$v=|V|^2$、流平衡不含损耗）；\fld{reconf\_loss\_mw} 是
$\sum r_{\text{pu}}\times\texttt{base\_mva}$（假设 1 pu 电流）的拓扑代理量，\textbf{非}潮流损耗，
需另跑完整潮流获取精确损耗。VSC 无功为近似。\texttt{ValidityFlags} 明确区分已解的图/线性拓扑
模型（\fld{ac\_lindistflow\_enforced}）与完整重构后交直流 PF/OPF 证书
（\fld{post\_power\_flow\_validated}、\fld{full\_hybrid\_opf\_validated}）。\fld{model\_scope}
报告 \texttt{"hybrid-acdc-topology-lindistflow"}（开启 PF）或
\texttt{"hybrid-acdc-topology-connectivity"}（仅连通性）。

\section{功能定位与架构}
\subsection{公共入口族}
\begin{center}\footnotesize
\begin{tabular}{@{}L{6.6cm}L{3.2cm}L{5.2cm}@{}}
\toprule
入口 & 定义位置 & 说明\\
\midrule
\srcpath{run\_topology\_reconfiguration(sys, opt)} & topology\_reconfiguration.hpp & 交直流混合故障后重构 MILP\\
\srcpath{solve\_optimal\_reconfiguration(ac\_sys, opt)} & topology\_analysis.hpp & 配电 ONR（LinDistFlow MILP + 验证潮流）\\
\srcpath{solve\_optimal\_reconfiguration(sys, opt)} & 同上 & \texttt{HybridPowerSystem} 重载（先投影为规范模型）\\
\srcpath{is\_connected / is\_radial / count\_islands(ac\_sys)} & 同上 & 快速拓扑查询\\
\bottomrule
\end{tabular}
\end{center}

\subsection{关键数据结构}
\compmeta{TopoReconfOptions}{include/hacdcpf/network\_reconfiguration/topology\_reconfiguration.hpp}
主要字段：\fld{line\_failures}、\fld{faulted\_branches}/\fld{switchable\_branches}
（\texttt{BranchRef}）、\fld{enable\_pf}/\fld{enable\_voltage}/\fld{enable\_thermal}、
\fld{split\_domain\_trees}/\fld{allow\_dc\_mesh}/\fld{loss\_aware}、
\fld{lambda\_switch}/\fld{lambda\_loss}/\fld{lambda\_shed}/\fld{lambda\_island}、
\fld{v\_min\_pu}/\fld{v\_max\_pu}、\fld{max\_time\_s}、\fld{mip\_gap}、\fld{solver}。

\compmeta{TopoReconfResult}{include/hacdcpf/network\_reconfiguration/topology\_reconfiguration.hpp}
主要字段：\fld{feasible}/\fld{optimal}、\fld{open\_branches}/\fld{closed\_branches}
（\texttt{BranchRef}）、\fld{switch\_operations}/\fld{ordered\_switch\_actions}、
\fld{reconf\_loss\_mw}、\fld{milp\_objective}/\fld{obj\_terms}、\fld{model\_scope}、
\fld{validity}（\texttt{ValidityFlags}）、\fld{solver\_backend}/\fld{solver\_status}。

\compmeta{ONRResult}{include/hacdcpf/network\_reconfiguration/topology\_analysis.hpp}
主要字段：\fld{open\_branch\_ids}/\fld{closed\_branch\_ids}、\fld{milp\_objective}、
\fld{estimated\_loss\_mw}、\fld{verification\_pf}、\fld{feasible}/\fld{optimal}、\fld{bc\_stats}。

\section{核心数学模型}
在规范化混合边集 $\mathcal{E}_{\text{hyb}}=\mathcal{E}_{\text{ac}}\cup\mathcal{E}_{\text{dc}}
\cup\mathcal{E}_{\text{vsc}}$ 上建立 MILP：支路状态二元量 $\beta_e$，单商品流连通性
（Dantzig--Johnson），生成树基数约束
\begin{equation}
\sum_{e}\beta_e + \sum_{g}\gamma_g = n_b,
\end{equation}
其中 $\gamma_g$ 为孤岛松弛、$n_b$ 为母线数。LinDistFlow 约束以 $v_i=|V_i|^2$ 表征电压，
大 M 电压降与热约束为
\begin{align}
\bigl\lvert v_j - v_i + 2 r_e P_e + 2 x_e Q_e \bigr\rvert &\le M\,(1-\beta_e), \\
\lvert P_e \rvert &\le P_e^{\max}\,\beta_e,
\end{align}
再叠加 VSC 传输与开关预算。目标为
\begin{equation}
\min\;\; \lambda_{sw}\sum_e \beta_e
       + \lambda_{loss}\sum_e r_e \beta_e
       + \lambda_{shed}\sum_i \bigl(s_i^{P}+s_i^{Q}\bigr)
       + \lambda_{isl}\sum_g \gamma_g .
\end{equation}
由原生分支定界（\srcpath{solver::solve\_milp\_bc}，含 MIR 割）或 HiGHS/SCIP 求解；辐射性按域
施加（VSC/DC-DC 换流器不入树，受 \fld{split\_domain\_trees}/\fld{allow\_dc\_mesh} 控制），
最优解后由 Newton 潮流交叉校核。损耗代理量为
\begin{equation}
P_{\text{loss}} \approx \sum_e r_{e,\text{pu}}\cdot S_{\text{base}} .
\end{equation}

\section{接口与参数}
\begin{paramtable}
\fld{v\_min\_pu} & $\underline{V}$ & pu & — & $0.95$ & 电压下限\\
\fld{v\_max\_pu} & $\overline{V}$ & pu & — & $1.05$ & 电压上限\\
\fld{lambda\_switch} & $\lambda_{sw}$ & — & $\ge 0$ & — & 开关动作权重\\
\fld{lambda\_loss} & $\lambda_{loss}$ & — & $\ge 0$ & — & 损耗（代理）权重\\
\fld{lambda\_shed} & $\lambda_{shed}$ & — & $\ge 0$ & — & 切负荷权重\\
\fld{lambda\_island} & $\lambda_{isl}$ & — & $\ge 0$ & — & 孤岛松弛惩罚权重\\
\fld{max\_time\_s} & — & s & — & $300$ & MILP 求解时限\\
\fld{mip\_gap} & — & — & $\ge 0$ & $0.01$ & MILP 相对间隙\\
\fld{enable\_pf} & — & 布尔 & — & \texttt{true} & 是否做最优解后验证潮流\\
\fld{solver} & — & — & — & \texttt{"auto"} & 求解器后端（auto/HiGHS/SCIP/原生 B\&C）\\
\fld{big\_m} & $M$ & — & — & $4.0$ & ONR 大 M 常数\\
\fld{default\_rate\_mva} & — & MVA & — & $10$ & ONR 默认支路容量\\
\end{paramtable}
交直流路径 \srcpath{run\_topology\_reconfiguration} 用 \texttt{BranchRef} 声明
\fld{faulted\_branches}/\fld{switchable\_branches}；ONR 路径 \texttt{ONROptions} 用
\fld{switchable\_branch\_ids} 声明可切支路，结果 \fld{verification\_pf} 承载验证潮流。

\section{验证与边界局限}
\subsection{验证与校核}
注册测试：\srcpath{test\_reconfig\_options}（\srcpath{run\_topology\_reconfiguration} 选项覆盖）、
\srcpath{test\_topology\_crossval}（两入口主交叉校核）、
\srcpath{test\_distribution\_pipeline}（流水线中的 ONR）、
\srcpath{test\_crossmodule\_integration}（\srcpath{solve\_optimal\_reconfiguration}）。

\subsection{边界与局限（如实标注）}
\begin{itemize}
  \item LinDistFlow 为线性近似（$v=|V|^2$、流平衡不含损耗）；\fld{reconf\_loss\_mw} 为
        $\sum r_{\text{pu}}\times\texttt{base\_mva}$ 拓扑代理量（假设 1 pu 电流），非潮流损耗，
        精确损耗须另跑完整潮流；
  \item VSC 无功为近似；\texttt{ValidityFlags} 区分 \fld{ac\_lindistflow\_enforced} 与
        \fld{post\_power\_flow\_validated}/\fld{full\_hybrid\_opf\_validated}；
  \item 旧式整数 ID 向量（\fld{open\_branch\_ids} 等）在 AC/DC/VSC 的 \fld{index} 重叠时含糊，
        须用 \texttt{BranchRef}；
  \item \srcpath{solve\_optimal\_reconfiguration(ac\_sys,\dots)}（ONR）为 AC-only LinDistFlow，
        混合交直流路径为 \srcpath{run\_topology\_reconfiguration}；\fld{model\_scope} 报告
        \texttt{"hybrid-acdc-topology-lindistflow"}（PF 开）或
        \texttt{"hybrid-acdc-topology-connectivity"}（仅连通性）。
\end{itemize}
\end{document}
